\begin{lemma}[Native regular injectivity in incident-bond coordinates]%
    \label{lem:peps_native_regular_incident_injectivity}
    \lean{TNLean.PEPS.torusIncidentCoordinatesEquiv,
        TNLean.PEPS.torusIncidentCoordinatesEquiv_apply,
        TNLean.PEPS.torusIncidentCoordinatesEquiv_smul,
        TNLean.PEPS.regularSiteMap_torusIncidentSite_eq_siteMap,
        TNLean.PEPS.IsGInjective.isGInjective_torusIncidentSite,
        TNLean.PEPS.torusGClosure_one_eq_stateCoeff_torusIncidentSite}
    \leanok
    \uses{lem:peps_torus_singleton_geometry, def:peps_g_injective}
    Let $G$ be a finite group and let a four-leg tensor be $G$-injective
    for the regular representation. On a native square torus with
    both periods at least three, identify its four virtual labels with
    the top, right, down, and left incident bonds. The resulting
    virtual-to-physical map is $G$-injective on the actual incident
    graph coordinates. Its ordinary closed graph contraction equals
    the native closure with both closure elements equal to the identity.
\end{lemma}

\begin{proof}\leanok
    The four actual legs are a bijection. The resulting bijection of
    virtual configurations intertwines simultaneous regular
    translation. Changing variables in the finite contraction sum
    identifies the two site maps. Invariance and injectivity on
    invariant vectors are therefore preserved. Identity closure
    elements insert identity operators on every seam bond, so the
    closed contraction agrees with the ordinary graph contraction.
\end{proof}

\begin{theorem}[Single-site support of a native regular injective tensor]%
    \label{thm:peps_native_regular_singleton_support}
    \lean{TNLean.PEPS.IsGInjective.range_regionReducedDensity_torus_singleton,
        TNLean.PEPS.IsGInjective.rank_regionReducedDensity_torus_singleton}
    \leanok
    \uses{lem:peps_native_regular_incident_injectivity,
        def:peps_region_reduced_density, def:peps_region_ground_space}
    Place one regular $G$-injective four-leg tensor at every site of a
    native square torus whose periods are at least three. For its
    untwisted closed vector, the physical density of every site $v$
    satisfies
    \[
        \operatorname{ran}\rho_{\{v\}}=\mathcal G_{\{v\}},
        \qquad \operatorname{rank}\rho_{\{v\}}=|G|^3.
    \]
\end{theorem}

\begin{proof}\leanok
    \uses{thm:peps_regular_torus_singleton_support,
        lem:peps_native_regular_incident_injectivity}
    Express native $G$-injectivity on the actual incident bonds and
    apply the single-site support and rank statements.
\end{proof}

\begin{theorem}[Local relations forced by equal native regular closures]%
    \label{thm:peps_native_regular_closure_local_rigidity}
    \lean{TNLean.PEPS.regionGroundSpace_torus_singleton_eq_of_torusGClosure_one_eq,
        TNLean.PEPS.card_group_eq_of_torusGClosure_one_eq}
    \leanok
    \uses{lem:peps_native_regular_incident_injectivity,
        def:peps_g_injective, def:peps_region_ground_space}
    Consider two native four-leg regular $G$-injective tensors with
    the same physical space on a square torus whose periods are at
    least three. If their untwisted closed coefficient vectors agree,
    their single-site ground spaces agree at every vertex. If the two
    descriptions instead use regular representations of different
    finite groups $G$ and $H$, equality of these untwisted closed
    vectors forces $|G|=|H|$.
    These are necessary local relations toward a Fundamental Theorem;
    neither a group isomorphism nor a product of individual bond
    gauges is asserted.
\end{theorem}

\begin{proof}\leanok
    \uses{thm:peps_regular_torus_local_rigidity,
        lem:peps_native_regular_incident_injectivity}
    Identity native closures are the ordinary graph contractions.
    The incident-bond descriptions remain $G$-injective, so the local
    ground-space and group-order rigidity statements apply.
\end{proof}
